\section{Discussion}
\label{sec:discussion}

% ---------------------------------------------------------------
\subsection{When Territorial Decomposition Helps}

The cross-context results reveal a clear pattern: decomposition via
spectral clustering provides the greatest benefit in instances where
(i) the global MILP is large enough that runtime matters, and
(ii) the candidate set is dense enough that individual subproblems
remain substantially smaller than the global problem.
The student area satisfies both conditions---it has nearly 5,000 demand
nodes, 67 existing stores, and 3,031 candidate blocks---and yields
a 67\% runtime reduction with less than 0.1 percentage-point coverage loss.
The city area has a smaller candidate set relative to its store density,
so the global instance is already fast; decomposition does not hurt
coverage but provides no speedup.
The industrial corridor violates condition (ii): with only 6 existing
stores covering 2,693 blocks, the vast majority of blocks qualify as
candidates, and each subproblem is proportionally harder than the global MILP.

This finding has a practical implication for deploying the framework
at scale: the benefit of decomposition should be estimated from the
ratio of candidate blocks to demand nodes before committing to the
clustered approach.
A simple rule of thumb is that decomposition pays off when the candidate
fraction exceeds roughly 55--60\% of all blocks in at least two thirds
of the resulting clusters; otherwise, the global solve is likely
faster and the decomposition gap is not worth accepting.

% ---------------------------------------------------------------
\subsection{Proportional Allocation and Coverage Equity}

The fixed-seven-per-cluster rule of the original formulation creates
a systematic inequity: clusters with larger demand receive the same
budget as clusters with smaller demand.
In the student area, this means the cluster with 1,282 blocks received
the same seven openings as the cluster with 212 blocks.
The demand-proportional rule (\cref{eq:proportional}) corrects this by
tying the budget to weighted demand share, resulting in allocations of
11, 11, 10, 5, 3, and 2 openings across the six clusters.
The coverage gain from proportionality is not dramatic in aggregate
(47.77\% vs.\ 44.51\% under fixed allocation in the student area),
but it substantially improves the per-capita coverage ratio in
the largest clusters, which is the operationally relevant measure
for a network expansion decision.

% ---------------------------------------------------------------
\subsection{The Role of MCA Weighting}

The MCA modifier shifts the optimization toward blocks with higher
socio-urban quality: better infrastructure, more complete urban services,
and higher educational attainment.
This reflects a deliberate commercial choice---the model is intended to
identify locations with the highest potential foot traffic, not merely
the highest population concentration.
The marginal coverage gain from MCA weighting (roughly 1--2 percentage
points in weighted coverage) is modest but directionally consistent
across all tested instances.

A more important contribution of MCA is its interpretive value:
the first MCA dimension provides a spatial visualization of socio-urban
quality that planners can inspect directly, and the boxplots by cluster
(\cref{fig:mca_by_cluster}) confirm that spectral territories differ
not only geometrically but also in their socioeconomic profile.
This means the cluster-level optimization is solving genuinely distinct
subproblems, not merely partitioning the same demand distribution.

The sensitivity analysis confirms that $\beta = 0.25$ is a stable
default: coverage and distance change by less than 1\% across
$\beta \in [0.1, 0.5]$, and population remains the dominant demand
signal (correlation above 0.97) throughout that range.

% ---------------------------------------------------------------
\subsection{Global vs.\ Clustered Optimality}

A reviewer concern was that cluster-wise optimization sacrifices global
optimality.
This concern is valid in principle: solving subproblems independently
may prevent a facility at the boundary of two clusters from serving
demand in both.
The empirical results, however, show that this effect is small in practice:
coverage losses range from 0.1 to 6.8 percentage points across the three
instances, and the weighted mean distance under the clustered model is
actually \emph{lower} than the global optimum in two of the three cases
(by 1.2~m and 0.4~m respectively), because the proportional allocation
concentrates openings in high-demand areas more effectively than the
single global problem under a fixed equal-split assumption.

The industrial corridor is the only case where the coverage gap is
commercially meaningful ($\approx$2.2 percentage points).
In that setting, a practitioner should consider running the global
\mbox{p-median} directly, since the clustered model provides neither
a runtime advantage nor a coverage benefit.

% ---------------------------------------------------------------
\subsection{Limitations}

Several limitations bound the scope of the reported results.
The framework optimizes coverage and access distance but does not
incorporate acquisition or operating costs, which can substantially
reorder candidate sites in practice.
The MCA weights depend on quintile discretization, so some within-quintile
variance is intentionally discarded.
The coverage rate is conditional on $D_{\max}$, and lower values
of the service radius reduce coverage monotonically, as shown
in the grid search; the reported 44--59\% coverage at $D_{\max} = 366$~m
is therefore a function of this deliberate constraint.
Finally, no external validation against observed sales, foot traffic,
or revenue data is possible here; the proposed model is validated against
the global \mbox{p-median} optimum as an internal benchmark, and the
question of whether MCA-weighted blocks do in fact generate higher
commercial activity remains an empirical question for future work.
